\documentclass{ltjsbook}
\usepackage{docmute} 
\usepackage{../myStyle} 

\begin{document} 
\chapter{Rをつかった因子分析と尺度作成法} 
\label{chapter10} 


    % \item[調査研究の手順] 心理尺度の作成研究の手続きを外観する。まず構成概念の設定，定義，妥当性を考えた上で，具体的な項目を選出し，テストデータを取る。探索的な因子分析によってその因子的妥当性を確認し，標準化のための本調査を行う。あるいは1次元性を確認した上で，IRTによって反応段階の確認，項目母数の確認，テスト情報関数の確認などが必要である。尺度の翻訳や検証的妥当性のチェックなどについては，構造方程式モデリングによる分析を行うのでここでは扱わず，参照するにとどめる。
    % \item[共通性推定の問題] 分析にあたって，改めて因子分析モデルの行列表現を提示し，行列の固有値分解によって因子負荷量が求められることを確認する。しかしその際，共通性をどのように推定するかの問題が残されていたことを確認し，そのためにいくつかの方法が提案されていることを理解する。これらは因子分析を行う上で，推定方法のオプション指定に関わってくる点であり，ソフトウェアが変わっても同様の指標が必要であることをみる。
     
    % \begin{flushright}
    %     $\to$\citet{kosugi2018}のPp.91--94.
    % \end{flushright} 
  
    % \item[fa関数と探索的因子分析] 探索的因子分析の手続きを\texttt{fa}関数を使いながら考える。探索的因子分析の場合は因子構造，因子負荷量について何ら前提を置かないため，因子数の推定から始めなければならない。まずは\texttt{fa.parallel}関数でスクリープロットを描画する。スクリープロットを読むときの形状について確認する。続いて因子数と共通性推定方法を定めた上で\texttt{fa}関数を実行し，因子負荷量や共通性などアウトプットを確認する。続いて解釈を簡単にするために因子軸の回転を行うことを解説し，実行のために\texttt{rotate}オプションを追加することをみる。回転前の結果との比較，また直交回転と斜交回転の違いを確認する。
 
     
    % \begin{flushright}
    %     $\to$\citet{kosugi2018}のPp.81--91.
    % \end{flushright} 
 
    % \item[因子得点の算出] 因子数と因子負荷量が明らかになると，そこから逆算的に因子得点を計算できる。\texttt{fa}関数には\texttt{scores}オプションをつけることで，出力されたオブジェクトから因子得点を取り出すことができるのをみる。こうした方法とは別に，項目同士の素点の平均から因子得点を計算することもある。これは推定値を実体とすることの懸念が出発点であり，その長所と短所を把握しておくことが必要である。この簡便法は平均値情報を含んでいるため，尺度カテゴリに依拠した解釈が可能である。また取り出された因子得点と簡便的因子得点の相関を見ることを確認する。
    % \item[因子分析の注意点] 因子分析を行う上で注意しなければならないのは，因子が実体としてあるのではなく，あくまでも準備された項目群の相関関係から得られる基底に過ぎないことを理解する点である。因子分析の流れの中では因子に命名することが1つの手順としてあるが，言葉として確定するとあたかもそれがあるかのように考えられてしまうこと，それしかないように考えられてしまうことの危険性を理解する。心理尺度の呪いやてっちゃんの手品になってしまわないように注意し，常に元の項目群に戻って考える必要があることをしっかりと理解する。
    



\section{課題}

\end{document} 
